\documentclass[10pt]{article}
\usepackage[a4paper,margin=2.2cm]{geometry}
\usepackage{amsmath,amssymb,booktabs,microtype,enumitem,array}
\usepackage[hidelinks]{hyperref}
\setlist{itemsep=2pt,topsep=3pt}
\setlength{\parskip}{4pt}\setlength{\parindent}{0pt}
\newcommand{\cfg}[1]{\texttt{#1}}

\title{RealPDE Track~2 (Long-Term Test-Time Adaptation):\\ the final submission in detail}
\author{Team nmucke}
\date{27 September 2026}

\begin{document}
\maketitle

\begin{abstract}
This document describes the final submission (Codabench submission 948041, archive \texttt{f11d2c1d}). It scored
82.452, first on the Track~2 development-phase leaderboard, and won Track~2 after the organisers' final evaluation,
in which they retrained the top methods. The model is a packed ensemble of three sim-pretrained, real-fine-tuned and weight-souped SimVP-style networks,
wrapped in a causal, gradient-free test-time pipeline. The pipeline has several parts: context normalisation,
a per-trajectory online least-squares correction of the ensemble head, calibration of the fluctuation energy,
damping of low-energy fluctuations, interval bounds chosen per trajectory to maximise the interval score, and a
detector that resets the running statistics after an abrupt regime change. The whole step runs as CUDA graph
replays scheduled over several GPU streams, using kernels that are bit-exact to the unfused computation. The
submission is rebuilt from the released data by one command.
\end{abstract}

\tableofcontents

%======================================================================
\section{Task and scoring}

Each hidden trajectory of real PIV measurements behind a NACA4418 airfoil (angle of attack
$0$--$20^\circ$, Reynolds number $3\,000$--$28\,000$) is revealed window by window. At step $t$ the model
receives the normalised input window $x_t\in\mathbb{R}^{1\times20\times32\times64\times3}$ (channels
$u,v,p$; $p$ is zero and unscored on real data) and the ground truth $y_{t-1}$ of the previous step. It returns
a prediction $\hat y_t$ and, optionally, per-element bounds $[\ell_t,h_t]$. The input and target windows
overlap in stride~20, so window $k$'s target is window $k{+}1$'s input. Everything inside
\cfg{ttt\_step} and \cfg{reset\_ttt\_state} is timed.

Five subscores are combined into the final score; the combination is not published. Our fit to the board is
$\approx 0.48\,\text{rel}_{L2} + 0.07\,\text{TKE} + 0.10\,\text{MVPE} + 0.09\,\text{time} + 0.28\,\text{SPS}$.
The Safe Prediction Score is computed per element as
\begin{equation}
  \text{elem}=\bigl[0.5(1-p_{\text{dm}})+0.3(1-p_{\text{tke}})+0.2(1-p_{\text{mvpe}})\bigr]\,
  e^{-(h-\ell)/\sigma_g}\,\mathbb{1}[\ell\le y\le h],\qquad p_x=\frac{e_x}{0.5+e_x},
\end{equation}
where $e_x$ are the per-window relative errors of the prediction. The per-window TKE error therefore enters
both the TKE subscore and the SPS accuracy factor. The time subscore is $100/(1+\sqrt{t/0.72896\,\text{s}})$.

%======================================================================
\section{Data}

Only the released \cfg{train\_sim} (100 simulation files) and \cfg{train\_real} (82 files) are used; no
external data or pretrained weights. \cfg{train\_real/7575\_0.h5} duplicates another file (organizer
announcement) and is excluded at every stage, which leaves 81 real files. Simulation data are subsampled to
the real $32\times64$ grid and rescaled to physical units, and $p$ is masked, so that pretraining sees exactly the
real input layout. The backbone uses no augmentation; the energy-net fits use an on-the-fly flip in $y$.

Development decisions used leave-conditions-out (LCO) folds that hold out whole angle-of-attack / Reynolds
regimes:
\begin{itemize}
\item \textbf{fold1:} used for selection.
\item \textbf{fold0:} the primary held-out fold.
\item \textbf{foldR:} held-out Reynolds numbers.
\item \textbf{Rx:} Reynolds extrapolation.
\item \textbf{foldS-lo, foldS-hi:} simulation-only regimes, as beyond-range surrogates.
\item \textbf{27 stress sets:} regime switches, velocity scaling, PIV dropout, noise, long trajectories and others.
\end{itemize}
The private test contains unseen parameter regimes, so every component below also had to be non-negative on
the private-regime proxies (foldR, Rx and the foldS sets).

%======================================================================
\section{Backbone and training}

\subsection{Architecture}
The backbone is a SimVP-style network with gated spatiotemporal-attention blocks
(\cfg{b5\_h192\_blk6\_rf11}): encoder width 128, hidden width 192, 6 gSTA blocks and receptive field 11,
3.39\,M parameters in total.

\subsection{Causal context normalisation}
The backbone is wrapped in causal context normalisation, which is the core test-time adaptation. For each
trajectory the wrapper keeps a per-pixel running mean and standard deviation of every frame revealed so far:
input frames and previously revealed targets, never the current target. These two maps are fed to the network
as extra channels, and the network's output is de-standardised with them. This adapts the operating point to
the trajectory's own regime with no gradient step. The state is cleared in \cfg{reset\_ttt\_state}.

\subsection{Loss}
The loss is a bounded Welsch (correntropy) loss, mixed 75/25 between physical and per-channel-standardised
space, with normalisers fixed from the first batch of each training stage.

\subsection{Training stages}
\begin{enumerate}
\item \textbf{Sim pretraining.} All 100 simulation files. AdamW, learning rate $5\cdot10^{-4}$, no weight decay,
  batch size 8, 3\,\% warmup and then cosine decay over a 20k-iteration schedule. The run is snapshotted at
  iteration 10k, a length selected on fold1.
\item \textbf{Real fine-tuning.} 8k iterations on the 81 real files with the same optimiser, starting from the
  sim snapshot.
\item \textbf{Weight soup.} Four fine-tunes (seeds 0--3) from the same snapshot are averaged uniformly in weight
  space (computed in fp64, cast to fp32). Because the members share one loss basin, the soup is better than any
  single member at no extra inference cost.
\item \textbf{Ensemble ($K=3$).} Stages 1--3 are repeated with sim-pretrain seeds 0, 1 and 2. Diversity comes
  from the pretrain seed; souping across seeds is invalid because the members are in different basins. The three
  soups are packed into one batched network (10.2\,M parameters): stem concatenation, grouped convolutions and
  batched $1{\times}1$ convolutions. The output is the members' mean in standardised space.
\item \textbf{fp16 core.} Everything inside the packed backbone runs in half precision; the surrounding
  pipeline stays fp32. The observed headroom to fp16 overflow is at least $1\,800\times$, and the composite
  changes by at most $|\Delta|\le0.0002$.
\end{enumerate}

\subsection{Fitted auxiliaries}
\begin{itemize}
\item \textbf{Prior:} a small per-condition seed that stabilises the first windows, with decaying influence.
\item \textbf{Energy net (``KE3''):} four energy nets (seeds 0--3) predict the fluctuation energy of the coming
  window from causal features. Their mean is written exactly as one dense network of $4\times$ width
  (block-diagonal weights; fp64 check $\le5\cdot10^{-9}$), so it runs as a single net on tensor cores.
\item \textbf{Interval-width net (``W2''):} a small CNN that predicts the log-ratio of the local error scale to the
  pipeline's running error scale. Its inputs are causal features, the lead index and the previous window's error
  map (EMA). It is fitted only on out-of-sample residuals from 4-group cross-fitting over the 81 files.
\end{itemize}

%======================================================================
\section{The online pipeline (per step, all causal)}

Every quantity below uses only $x_t$, the revealed $y_{<t}$ and fixed constants. Grids, thresholds and
constants are literals in \cfg{submission.py}; the retrain chain fits none of them. The stages run in the order
listed.

\subsection{Regime-change detector (REC)}
The detector computes $d=|\log(\bar U_{\text{win}}/\bar U_{\text{run}})|$ from the mean physical $u$ of the
current input frames and the running mean since the last reset. If $d>\log1.5$, the context-normalisation
accumulator, the error, centre and blend moments, and the least-squares moments of
Section~\ref{sec:lsk} restart. Mean-type statistics are rescaled by
$\bar U_{\text{win}}/\bar U_{\text{run}}$ and given a pseudo-count of 6.

On all 81 real files the statistic stays below 0.085, so the detector never fires and REC is bit-identical to
its absence on normal streams. On synthetic regime-switch streams it recovers within 0--2 windows.

\subsection{Online least-squares correction of the ensemble head (LSK)}\label{sec:lsk}
The packed head computes $\frac1K\sum_k (W_kh_k+b_k)=W_{\text{cat}}H$ with
$W_{\text{cat}}=[W_1\dots W_K]/K$. Let $V$ hold the top-8 right singular vectors of the $u,v$ rows of
$W_{\text{cat}}$; they are computed at model construction from the trained weights, so no extra fit is needed.
The feature vector is $\phi=[V^\top H;1]$, and $s$ is the causal context standard deviation. For each channel,
from the revealed error $e=y_{t-1}-\hat y^{\text{raw}}_{t-1}$:
\begin{align}
  \Phi&\leftarrow\rho\Phi+\textstyle\sum s^2\phi\phi^\top, &
  B&\leftarrow\rho B+e\,(s\phi)^\top, &
  M&=B(\Phi+\lambda I)^{-1},
\end{align}
with $\rho=0.97$ and $\lambda=30\,\mathrm{tr}(\Phi_1)/9$ (fixed from the first revealed window). The raw $u,v$
output of all 20 leads is then corrected by $(M\phi)s$ before any downstream stage.

The moments are held in fp64 and the $9\times9$ solve uses \cfg{solve\_ex} with no host synchronisation. A
non-finite guard skips bad pairs. The moment update runs on a side GPU stream, so only one GEMM and the apply
sit on the step's critical path. A build-time self-test uses the configured $\rho$ and falls back to ``LSK off''
on any mismatch. REC multiplies the moments by its keep factor.

\subsection{Energy calibration}
The predicted fluctuation amplitude is rescaled so that its energy matches a blend of the energy-net prediction
and the running mean of revealed energy. The rescale adds the mean back so that variance is preserved, and the
gain is clamped to $[1/20,\,20]$.

The blend weight $b\in\{0.45,0.60,0.75,0.85,0.925\}$ is chosen per trajectory. Every candidate is computed each
step, batched, and the one with the best EMA ($\rho=0.95$) of the previous window's official composite is
emitted. Before any evidence exists, $b=0.85$.

\subsection{Prior and advection blend}
The prediction is blended with an advected copy of the most recently revealed frame. The ridge strength of this
blend's least-squares fit is chosen per trajectory from $\{0.125,0.25,0.5,1\}$ by the same previous-window-score
rule.

\subsection{Low-energy fluctuation damping (LOWKE)}
Let $m_t$ be the 20-frame mean of the prediction and $K_p=\tfrac12(\mathrm{var}_t u+\mathrm{var}_t v)$ its
energy map. For windows $\ge2$,
\begin{equation}
  g=\mathrm{clamp}\!\Bigl(\bigl(K_p/(1.5\,\overline{K_p})\bigr)^{0.25},0,1\Bigr),\qquad
  \hat y\leftarrow m_t+g\,(\hat y-m_t).
\end{equation}
This removes weakly correlated fluctuations where the model itself predicts calm flow. rel$_{L2}$ and SPS rise
and TKE drops slightly. The bounds and internal state are untouched.

\subsection{Interval bounds}
\begin{itemize}
\item \textbf{Width:} the W2 half-width, multiplied by the running error scale and by an adaptive multiplier
  $\kappa\in\{0.85,0.925,1,1.075,1.15\}$. $\kappa$ is chosen separately for leads 1--2 and leads 3--20 by the EMA
  of the exact SPS element score on the previous window.
\item \textbf{Centre:} the centre is shifted toward the raw ensemble output,
  $c'=c+w(\hat y^{\text{raw}}-c)$. The weight $w\in\{0.30,0.45,0.60,0.75,0.90\}$ is chosen per trajectory by the
  same rule. The fallback is $w=0.75$, and window~0 uses $w=1$.
\item \textbf{Scored mask:} the per-lead error maps use the scorer's exact scored mask,
  $|y|>10^{-6}$ on the normalised target.
\end{itemize}
Bounds are returned on every step of every run, as the all-or-nothing rule requires.

\subsection{Causality summary}
\begin{center}\small
\begin{tabular}{@{}p{5.2cm}p{9.8cm}@{}}\toprule
Stage & Information used \\\midrule
Context normalisation, REC & frames revealed up to $t-1$; the current step's input frames\\
LSK & revealed errors up to $t-1$; the current step's own head features\\
Energy / blend / ridge / $\kappa$ / $w$ selectors & scores of previous windows only (EMA)\\
Prior, W2, energy nets & fixed after training (training files only)\\
Current target $y_t$ & never used for step $t$\\\bottomrule
\end{tabular}
\end{center}

%======================================================================
\section{Execution}

\subsection{CUDA graphs and fallbacks}
From window 2 onward the entire step is one CUDA-graph replay. Windows 0 and 1 have their own captured graphs,
and every graph has a bit-exact eager fallback. Each graph is self-tested against eager execution at model
construction, and any failure degrades gracefully. An RNG probe guards against capture-time corruption of the
CUDA random state.

\subsection{Kernel diets}
All diets preserve outputs exactly, except KD5, which changes rounding order only
($|\Delta\text{composite}|\le0.0004$).
\begin{itemize}
\item \textbf{Fused selectors:} the candidate scoring is fused into fewer kernels.
\item \textbf{Leaner step body:} the backbone and ``mega'' stage are rewritten with fewer kernels.
\item \textbf{KD5:} near-exact convolution rewrites.
\item \textbf{First-two-window graphs:} windows 0 and 1 are captured with a side stream and private memory pools.
\end{itemize}

\subsection{Multi-stream schedule}
The step graph is captured with fork/join edges. Work that does not depend on the backbone output (error and
loss bookkeeping, running variances, width-net updates, adaptive selectors, input-only preparation) runs on
parallel streams alongside the backbone. The LSK update and the LOWKE gain also have their own streams. Every
cross-stream read-after-write and write-after-read hazard is ordered by an event.

This was verified in three ways: a race checker using transitive happens-before (0 conflicts); replays with
deliberate delays that invert each fork/join completion order; and runs under foreign GPU load and different
connection limits. All comparisons matched bit for bit.

\subsection{Fused elementwise kernels}
The longest elementwise chains on the critical path are fused into \cfg{torch.cuda.jiterator} kernels. They are
written with round-to-nearest intrinsics (\cfg{\_\_fmul\_rn}, \cfg{\_\_fadd\_rn}, \dots) so that NVRTC cannot
contract them into fused multiply-adds, which keeps them bit-identical to the unfused graph. They are compiled
and bit-tested at construction, outside the timed region; any failure falls back.

\subsection{Transport across the harness process boundary}
The official harness runs the model in a worker subprocess.
\begin{itemize}
\item \textbf{By-value outputs (BVR):} outputs are returned as a zero-copy tensor subclass whose pickled form is
  its values, so the parent receives ordinary tensors with identical values. This proved decisive on the board.
\item \textbf{Other transport paths:} fd-shared single-storage outputs, a listener warm-up during construction
  (untimed), and direct CUDA I/O cover the other possible harness designs. They are inert where they do not
  apply.
\end{itemize}

\subsection{Resulting speed}
The eval harness reports 4.84\,ms per step. Of that, about 2.6\,ms is our GPU work on the A800-class device;
the rest is harness transport.

%======================================================================
\section{Validation evidence for each component}

Each component was adopted under a gate declared in advance: selection on fold1, then a single held-out
evaluation on fold0, foldR, Rx and the foldS sets. Most were also reviewed by an independent Fable reviewer and
verified on the exact archive: bit-identity checks, out-of-process harness emulation, fuzzing, reset probes and
determinism.

The table lists the change in the local composite (time held fixed) at the moment each component was added to
the stack of that time. It therefore shows increments, not a clean ablation of the final stack. Board transfer
is about $0.6\times$ fold0.

\begin{center}\small
\begin{tabular}{@{}lrrrrrr@{}}\toprule
Component & fold1 & fold0 & foldR & Rx & foldS-lo & foldS-hi\\\midrule
Sim pretraining            & +0.105 & +0.184 & +0.064 & +0.177 & +0.186 & $-$0.181\\
Soup of 4                  & +0.022 & +0.075 & +0.048 & +0.041 & +0.024 & $-$0.001\\
2nd ensemble member        & +0.080 & +0.091 & +0.082$^a$ & +0.060$^a$ & +0.073$^a$ & +0.084$^a$\\
3rd member + fp16 core     & +0.048 & +0.039 & -- & -- & +0.026 & $-$0.015\\
Centre blend               & +0.110 & +0.089 & +0.114 & -- & -- & --\\
Adaptive centre weight     & +0.001 & +0.001 & -- & -- & +0.002 & +0.055\\
Adaptive width / energy blend & +0.017 & +0.027 & +0.011 & -- & +0.024 & +0.046\\
Gain clamp 20              & +0.006 & +0.008 & +0.031 & +0.019 & $-$0.004 & +0.000\\
Exact scored mask          & +0.003 & +0.005 & +0.004 & +0.009 & $-$0.001 & +0.005\\
Beyond-range hedge (AXE)   & +0.000 & +0.009 & -- & -- & +0.009 & +0.235\\
4-seed energy net (KE3)    & +0.018 & +0.043 & +0.020 & -- & +0.057 & +0.115\\
LOWKE (net of time)        & +0.022 & +0.033 & +0.021 & -- & +0.037 & +0.035\\
LSK ($\lambda_{\text{rel}}=30$)& +0.018 & +0.041 & +0.027 & +0.034 & -- & --\\
LSK $\rho$ 0.9$\to$0.97$^b$ & +0.004 & +0.010 & +0.005 & +0.007 & +0.008 & +0.033\\\bottomrule
\end{tabular}
\end{center}
{\footnotesize $^a$ measured on an independent 4090 build.
$^b$ Missed the fold0 bar of +0.015 declared in advance. It was adopted as a disclosed override because it was
non-negative on every held-out and private-proxy set, and because an independent refit on the 4090 reproduced
it (fold0 +0.009, 5/5 sets $\ge0$).}

Notes on the table:
\begin{itemize}
\item REC is bit-identical on every normal stream and gains +17--22 composite on a regime-switch file.
\item Resetting the LS moments on REC is bit-identical off-switch and gains +0.036 on the switch set.
\item The timing layers are output-identical except KD5 (rounding order only). On a 3090 they save:
  fused selectors $-0.11$\,ms, leaner step body $-0.11$\,ms, KD5 $-0.115$\,ms, first-two-window graphs
  $-0.05$\,ms, multi-stream schedule $-0.42$\,ms ($-0.50$\,ms on an A100), fused kernels $-0.12$\,ms and the
  final critical-path cuts $-0.07$\,ms.
\item By-value transport saves about 0.7\,ms per step on the board.
\end{itemize}

In a 27-set stress battery with unseen-regime surrogates, the pipeline beat the best trivial predictor on every
set, with no exceptions and no non-finite outputs.

%======================================================================
\section{Result}

\begin{center}
\begin{tabular}{lcccccc}\toprule
& rel\_l2 & TKE & MVPE & time & SPS & final\\\midrule
Development leaderboard, 27 Sep 2026 & 94.939 & 81.017 & 94.500 & 92.464 & 45.211 & \textbf{82.452}\\\bottomrule
\end{tabular}
\end{center}
The pre-registered prediction was $82.455$ ($82.40$--$82.51$). The time subscore corresponds to 4.84\,ms per step.
This was first on the development-phase leaderboard (Codabench submission 948041). In the organisers' final
evaluation, for which they retrained the top methods, this submission won Track~2.

%======================================================================
\section{Reproducibility and environment}

\begin{itemize}
\item \textbf{One command:} \cfg{bash scripts/retrain\_all.sh}, run in \cfg{final\_submission/realpde/} of the
  code release, rebuilds \cfg{submission.zip} from the released
  data. The chain has 13 idempotent steps and uses seed 0 throughout.
\item \textbf{Wall time:} about 2--3\,h on an RTX~3090 (2\,h\,47 measured from a fresh clone), about 1\,h on an RTX~4090 and about 2\,h\,15 on an A100 with the evaluation software
  versions. A from-scratch A100 retrain scored within $\pm0.03$ of our build on every local subscore.
\item \textbf{What is reproduced exactly:} code entries are byte-identical, and on the same GPU so are
  \cfg{model.pth} and \cfg{prior.npz}. Under GPU contention the sim-pretraining snapshot can differ bitwise, so
  the correctness bar is score parity, not byte identity.
\item \textbf{Runtime:} Python~3.10, PyTorch~2.2.2 (CUDA 12.1), numpy~1.26. The NVRTC shipped with torch is
  needed for the fused kernels, which fall back if it is missing. Inference is offline and reads nothing from disk
  beyond the archive; the extracted archive is about 60\,MB.
\item \textbf{Robustness of the chain:} an opt-in serial mode supports clusters that allow one process per GPU.
  Reused artefacts are validated before a step is skipped. An open-file audit confirmed that only the released
  training data are read.
\end{itemize}

%======================================================================
\section{Limitations}
\begin{itemize}
\item TKE and the remaining SPS gap are limited by the backbone's forecast correlation at late leads and high
  Reynolds numbers. Output-side corrections tested on top of the final stack (spectral, temporal and physical
  projections; phase warps; per-trajectory reduced models) did not improve it.
\item The highest-Reynolds 0$^\circ$ files remain the weakest regime: LSK is slightly negative on a few of
  them.
\item Synthetic bursts of all-zero frames degrade the LS correction. No released real frame has this property.
\item Timing levers that depend on the harness were validated on emulations and on the board; the exact private
  harness is not public.
\end{itemize}

\end{document}
